\centerline{\bf \Huge Числовые ребусы II}

\begin{flushright}
        \scriptsize{Добавлю вашу цитату}
\end{flushright}

\textbf{В ребусах с участием букв действует следующее правило: }

\textit{Одинаковым буквам должны соответствовать одинаковые цифры, разным буквам — разные цифры.}

\textbf{Решить ребус значит найти ВСЕ решения и доказать, что других нет!}

\problem{\textbf{1}} ТРАССА + ТРАССА = КОСМОС

\problem{\textbf{2}} КНИГА + КНИГА + КНИГА = НАУКА

\problem{\textbf{3}} БАРАН + БАРАН + БАРАН + БАРАН = СТАДО

\problem{\textbf{4}} АИСТ + АИСТ + АИСТ + АИСТ + АИСТ = СТАЯ

\problem{\textbf{5}} FORTY + TEN + TEN = SIXTY

\problem{\textbf{6}} В + ЛЕС + ИЗ + НОРЫ = 3330

\problem{\textbf{7}} Л + $\dfrac{O}{M}$ + О + Н + $\dfrac{O}{C}$ = ОВ (Обе дроби правильные, то есть числитель меньше знаменателя)